\documentclass[11pt,a4paper]{article}
\usepackage[margin=2.2cm]{geometry}
\usepackage{ctex}
\usepackage{amsmath,amssymb,amsthm}
\usepackage{enumitem}
\usepackage{tasks}
\usepackage{array,tabularx}
\usepackage{multicol}
\usepackage{fancyhdr}
\usepackage{tikz}
\usepackage{graphicx}
\usepackage{xcolor}
\usepackage{needspace}

\pagestyle{fancy}
\fancyhf{}
\fancyhead[L]{\small 大模型高性能计算 期末试卷}
\fancyhead[R]{\small 第 \thepage\ 页}
\renewcommand{\headrulewidth}{0.4pt}

\setlength{\parindent}{0pt}
\setlength{\parskip}{0pt}

\newcommand{\blank}[1]{\underline{\hspace{#1}}}
\newcommand{\fn}[2]{\dfrac{#1}{#2}}

\newcommand{\figplaceholder}[2]{%
  \begin{center}
  \fcolorbox{gray!60}{gray!5}{%
    \begin{minipage}{0.88\linewidth}
    \vspace{0.3cm}
    \centering
    {\small\bfseries [配图占位符]}\\[0.2cm]
    {\small #1}\\[0.15cm]
    {\footnotesize\kaishu 图源：#2}
    \vspace{0.3cm}
    \end{minipage}}
  \end{center}
  \vspace{0.3cm}
}

\newcommand{\examsection}[2]{%
  \vspace{0.6cm}
  \noindent{\CJKfamily{SimHei}\bfseries\large #1\quad (#2)}\par
  \vspace{0.3cm}
}

\newcommand{\qitem}[2]{%
  \noindent\textbf{#1.}\;#2\par
  \vspace{0.3em}
}

\newcommand{\fourchoices}[4]{%
  \begin{enumerate}[label=(\Alph*),leftmargin=2.2em,itemsep=0pt,parsep=0pt,topsep=0pt]
    \setlength{\itemindent}{0pt}
    \item #1
    \item #2
    \item #3
    \item #4
  \end{enumerate}
  \vspace{0.4em}
}

\newcommand{\fourchoicestwo}[4]{%
  \noindent
  \begin{minipage}[t]{0.5\linewidth}\makebox[2.2em][l]{(A)}#1\end{minipage}%
  \begin{minipage}[t]{0.5\linewidth}\makebox[2.2em][l]{(B)}#2\end{minipage}\par
  \noindent
  \begin{minipage}[t]{0.5\linewidth}\makebox[2.2em][l]{(C)}#3\end{minipage}%
  \begin{minipage}[t]{0.5\linewidth}\makebox[2.2em][l]{(D)}#4\end{minipage}\par
  \vspace{0.4em}
}

\newcommand{\answerlines}[1]{%
  \vspace{0.2cm}
  \foreach \i in {1,...,#1} {
    \noindent{\color{white}\rule{\linewidth}{0.4pt}}\par\vspace{0.5cm}
  }
}

\begin{document}

\begin{center}
{\zihao{2}\CJKfamily{SimHei}\bfseries 大模型高性能计算 期末试卷}\\[0.3cm]
{\zihao{4}\kaishu 满分 150 分 \quad 考试时间 150 分钟}
\end{center}

\vspace{1.5cm}
\noindent\rule{\linewidth}{1pt}

\vspace{0.2cm}
\noindent{\CJKfamily{SimHei}\bfseries 注意事项：}
\begin{enumerate}[leftmargin=2em,itemsep=0pt,parsep=0pt,topsep=0pt]
  \item 本试卷共 \textbf{三} 大题，包括选择题、填空题、解答题，共 \textbf{24} 小题。
  \item 请将答案写在答题区指定位置；解答题须写出必要的文字说明、推导过程或演算步骤。
  \item 涉及数学公式推导时，请清晰写出关键中间步骤；只写最终结果不得分。
  \item 本试卷中如无特殊说明，GEMM 维度记为 $C_{M\times N}=A_{M\times K}B_{K\times N}$；HBM 峰值带宽记 $B_w$，单位 GB/s；峰值算力记 $P_{\mathrm{peak}}$，单位 TFLOPS；SRAM 容量记 $M_{\mathrm{sram}}$；Warp 大小默认 32；fp16 每标量 2 字节，fp32 每标量 4 字节。GPU 硬件参数默认指 NVIDIA H100-80GB，SXM 形态。
\end{enumerate}

\vspace{0.2cm}
\noindent\rule{\linewidth}{0.6pt}

\examsection{一、选择题}{本大题共 12 小题，每小题 5 分，共 60 分。在每小题给出的四个选项中，只有一项是符合题目要求的。}

\qitem{1}{现代 NVIDIA GPU，如 H100，将大量计算单元组织为流多处理器 SM。每个 SM 内含 CUDA Core 负责通用浮点与整数运算、Tensor Core 负责矩阵乘累加 MMA、Warp 调度器、片上可编程缓存 Shared Memory，以及线程私有的寄存器文件 Register File。线程以 Warp 为单位调度，默认每 Warp 32 线程，同一 Warp 内线程执行相同指令流，即 SIMT 模型。关于上述结构，下列说法\textbf{正确}的是}{}
\fourchoices{每个 SM 仅包含一个 CUDA Core，所有线程在同一时刻只能串行执行一条指令}{每个 SM 拥有独立的 Warp 调度器，可将 32 个线程组成一个 Warp 进行 SIMT 调度；若遇分支发散，Warp 内部分线程将空闲等待}{SM 内的 Shared Memory 与 L2 Cache 在物理上是同一块存储，只是软件命名不同}{寄存器文件位于 HBM 上，所有 SM 通过 PCIe 总线共享访问，延迟与 Global Memory 相同}

\qitem{2}{GPU 存储层级从快到慢依次为 Register、Shared Memory、L2 Cache、HBM 高带宽内存。各级存储在容量、带宽和访问延迟上差异显著：Register 最快但每线程仅有数千字节；Shared Memory 在 SM 内共享、需程序员显式管理；L2 为芯片级统一缓存；HBM 容量最大，H100 上达 80GB，但延迟最高。高效 Kernel 设计的核心目标之一，是将数据尽可能长时间停留在高速存储层。关于带宽与延迟，下列从左到右\textbf{带宽递减、延迟递增}的排序\textbf{正确}的是}{}
\fourchoices{Register $>$ Shared Memory $>$ L2 $>$ HBM}{HBM $>$ L2 $>$ Shared Memory $>$ Register}{Shared Memory $>$ Register $>$ HBM $>$ L2}{L2 $>$ HBM $>$ Register $>$ Shared Memory}

\qitem{3}{Roofline 模型用算术强度 $I$ 判断算子瓶颈，$I$ 的单位为 FLOP/Byte，表示每从 HBM 搬运 1 字节数据能完成多少浮点运算。硬件 Ridge Point 定义为 $I^*=\dfrac{P_{\mathrm{peak}}}{B_w}$，其中 $P_{\mathrm{peak}}$ 为峰值算力，$B_w$ 为 HBM 峰值带宽。实际吞吐近似为 $\min\!\left(P_{\mathrm{peak}},\, I\cdot B_w\right)$：当 $I<I^*$ 时算力无法跑满，性能受带宽限制；当 $I>I^*$ 时带宽足够，性能受算力限制。大模型 Decode 阶段 GEMM 的批大小维 $M$ 往往很小，算术强度通常极低。当 $I < I^*$ 时，该算子处于}{}
\fourchoicestwo{Compute-bound，计算受限，应优先提升 Tensor Core 利用率}{Memory-bound，带宽受限，应优先减少 HBM 访问量或提高数据复用}{Communication-bound，通信受限，应优先优化 NCCL 拓扑}{Launch-bound，启动开销受限，应优先减少 Kernel 数量}

\qitem{4}{矩阵乘法 $C_{M\times N}=A_{M\times K}B_{K\times N}$ 是大模型最核心的算子，广泛用于 Attention 投影与 FFN 线性层。每个输出元素 $C_{ij}=\sum_{k=1}^{K}A_{ik}B_{kj}$ 需要 $K$ 次乘加 MAC。在 FLOPs 统计中，1 次 MAC 计为 2 次浮点运算，即 1 次乘法加 1 次加法。因此，完成整个 GEMM 所需的浮点运算量为}{}
\fourchoicestwo{$MNK$}{$2MNK$}{$3MNK$}{$MN+NK+MK$}

\qitem{5}{Google TPU 与 NVIDIA GPU 代表两种不同的 AI 加速器设计哲学。GPU 采用 SIMT 单指令多线程模型，擅长通用并行计算和复杂控制流；TPU 采用专用二维脉动阵列 Systolic Array，权重矩阵预加载到阵列中，激活值按固定节奏流过阵列完成 MAC，数据流高度规整。TPU 在规整的大矩阵乘法上效率极高，但对不规则访存和控制流支持较弱。关于二者架构差异，下列说法\textbf{正确}的是}{}
\fourchoices{TPU 采用与 GPU 完全相同的 SIMT 模型，每个 TPU Core 可独立执行不同分支指令}{TPU 核心为二维脉动阵列，数据在阵列中按固定节奏流动完成矩阵乘累加，无需像 GPU 那样为每个线程维护独立 PC}{GPU 的 Tensor Core 与 TPU 脉动阵列在硬件数据流上完全相同，只是命名不同}{NPU 是专有名词，仅指华为昇腾芯片，与 TPU、GPU 在架构上无任何可比性}

\qitem{6}{算子融合 Kernel Fusion 是编译器与手写 Kernel 中的常用优化：将多个连续算子，如 RMSNorm、Linear、SiLU，合并为单个 CUDA 或 Triton Kernel，使中间结果保留在寄存器或 Shared Memory 中，避免写回 HBM 再读取。融合还能减少 Kernel Launch 开销，每次 Launch 约有 5--20$\mu$s 的 CPU 调度延迟。vLLM、TensorRT-LLM 等推理框架均大量采用融合。关于算子融合的收益，下列说法\textbf{错误}的是}{}
\fourchoices{融合多个算子可减少 Kernel Launch 次数，降低 CPU 侧调度开销，对短算子链尤其明显}{融合后中间激活可保留在片上存储，减少 HBM 读写往返}{算子融合一定能将任意两个算子的总执行时间精确减半，且不受寄存器数量和算子形状限制}{RMSNorm + Linear、SiLU + Mul 即 SwiGLU 门控，是工业界常见且收益明显的融合模式}

\qitem{7}{W8A8 量化推理将权重和激活均量化为 INT8，在 INT32 累加器上完成 GEMM 后再反量化到 fp16 或 fp32。设激活 $\mathbf{X}$ 使用 per-tensor 缩放因子 $s_x$，即整个张量共用一个 scale；权重 $\mathbf{W}$ 使用 per-channel 缩放因子 $s_w^{(j)}$，即每个输出通道 $j$ 独立 scale，精度更高。量化值为 $\mathbf{X}_{\mathrm{int8}}=\mathrm{round}(\mathbf{X}/s_x)$，$\mathbf{W}_{\mathrm{int8}}=\mathrm{round}(\mathbf{W}_{:,j}/s_w^{(j)})$。则等价的 fp32 线性层输出在工程实现中通常写为}{}
\fourchoices{$\mathbf{Y} = \mathbf{X}_{\mathrm{int8}} \mathbf{W}_{\mathrm{int8}}$，直接输出 INT8 结果，无需反量化}{先执行 INT32 累加 $\mathbf{Y}_{\mathrm{int32}}=\mathbf{X}_{\mathrm{int8}}\mathbf{W}_{\mathrm{int8}}$，再对每个输出通道 $j$ 乘 $s_x\cdot s_w^{(j)}$ 反量化到 fp16 或 fp32}{W8A8 严格要求权重和激活均使用 per-tensor 量化，不允许 per-channel}{INT8 GEMM 的累加必须在 INT8 域完成并直接输出，使用 INT32 累加器属于数值错误}

\qitem{8}{NVIDIA Tensor Core 自 Volta 架构引入，专门加速矩阵乘累加 MMA。以 FP16 输入为例，标准 MMA 指令形状为 $m16n8k16$：每次处理 $16\times8$ 的输出 tile，沿 $K$ 维每次累加 16 个元素；输入为 FP16 或 BF16，累加器为 FP32 以保证精度。CUTLASS、cuBLASLt 等库通过多级 Tiling 将大 GEMM 分解为适配 MMA 形状的子块。关于 Tensor Core，下列说法\textbf{正确}的是}{}
\fourchoices{Tensor Core 仅支持 FP32 输入与 FP32 累加，与 CUDA Core 无区别}{标准 FP16 MMA 指令形状为 $m16n8k16$，输入 FP16 或 BF16、累加器 FP32；高效 GEMM 仍需配合多级 Tiling}{Tensor Core 与 CUDA Core 互斥，同一次 Kernel Launch 中不能混用二者}{调用 cuBLAS 的 GEMM 会自动使用 Tensor Core，开发者无需关心 Tiling 与数据布局}

\qitem{9}{GPTQ、AWQ、SmoothQuant 是大模型 PTQ 训练后量化的三类代表性方法，实现逻辑各不相同。GPTQ 利用 Hessian 二阶信息逐列量化权重并传播误差；AWQ 识别对输出影响大的显著权重通道并按比例缩放；SmoothQuant 通过对角矩阵 $S$ 将激活值中的离群点平滑迁移到权重侧，使激活更易量化。下列关于三者实现逻辑的匹配，\textbf{正确}的是}{}
\fourchoices{GPTQ 基于 Hessian 二阶信息逐列量化权重；AWQ 基于显著权重通道缩放；SmoothQuant 通过 $S$ 矩阵将激活离群点迁移到权重侧}{GPTQ 基于激活值平滑迁移；AWQ 基于 Hessian 逐列量化；SmoothQuant 基于显著通道缩放}{三者均属于 QAT 量化感知训练，训练时需插入 FakeQuant 算子并完整反向传播}{三者均只量化权重，推理时激活值始终保持 fp32 不做任何压缩}

\qitem{10}{Triton 是面向 GPU 的高级 DSL，以 Block 为基本编程抽象：开发者描述每个 Block 做什么，编译器自动生成 Thread 与 Warp 索引、内存合并访问 Coalesced Access 和 Shared Memory 布局。相比之下，手写 CUDA 需显式管理 \texttt{threadIdx}、\texttt{blockIdx} 以及 Warp 级原语如 \texttt{\_\_shfl\_sync} 等底层细节。Triton 已广泛用于 FlashAttention、量化 GEMM 等自定义 Kernel。关于二者差异，下列说法\textbf{正确}的是}{}
\fourchoices{Triton 以 Thread 为基本编程单元，开发者需像 CUDA 一样手动计算每个线程的全局索引}{Triton 以 Block 为基本抽象，编译器自动处理 Thread 与 Warp 级并行及 Global Memory 合并访问}{Triton 无法编写自定义 GEMM Kernel，所有矩阵乘法必须调用 cuBLAS}{手写 CUDA 完全无法做算子融合，融合优化只能由 Triton 实现}

\qitem{11}{FlashAttention 通过 Tiling 将 $Q,K,V$ 分块加载到 SRAM，在片上完成 Online Softmax 并增量累加输出，避免将完整的 $N\times N$ 注意力矩阵 $S,P$ 写回 HBM。FlashAttention-2 进一步在 Warp 级别并行化、减少非矩阵乘 FLOPs，并优化了 SRAM 中的工作划分。关于 FlashAttention-2 的改进，下列\textbf{不包括}在内的是}{}
\fourchoices{减少非矩阵乘法的 FLOPs，如对 running statistic 的 rescale 操作进行合并与简化}{在 Warp 级别更好地并行化 attention 计算，提升 SM 利用率}{将 $Q$、$K$、$V$ 的 SRAM 分块策略优化为跨 Warp 分工协作}{将注意力计算的 HBM IO 复杂度从 $O(N^2d^2/M)$ 提升到 $O(N^3)$，以换取更低的 SRAM 占用}

\qitem{12}{XLA 和 TVM 是两类主流深度学习编译器。XLA 将 TensorFlow 或 JAX 计算图降低为 HLO 中间表示，执行算子融合 Fusion、公共子表达式消除 CSE、内存布局优化 Layout 等 Pass 后生成 LLVM 或 CUDA 代码。TVM 通过 AutoTVM 或 Ansor 对 Kernel 的 tile size、unroll factor、并行策略等超参进行自动搜索，在多种硬件后端上寻找最优实现。关于编译器优化，下列说法\textbf{正确}的是}{}
\fourchoices{XLA 的 HLO 图优化不支持算子融合，PyTorch Eager 模式下执行的算子无法被 XLA 优化}{TVM 通过 AutoTVM 或 Ansor 对 Kernel 的 tile size、unroll factor 等超参进行自动搜索，可面向 GPU、CPU 等多种后端}{TVM 是 Google 内部专用工具，只能为 TPU 生成代码，不支持 NVIDIA GPU}{XLA 的 Layout Optimization 与 Kernel Fusion 完全独立、互不影响，融合不会改变张量内存布局}

\examsection{二、填空题}{本大题共 6 小题，每小题 5 分，共 30 分。把答案填在题中横线上。}

\qitem{13}{Roofline 模型中，Ridge Point $I^*=P_{\mathrm{peak}}/B_w$ 是判断算子瓶颈的关键阈值。设 H100-80GB 的 FP16 Tensor Core 稠密峰值算力 $P_{\mathrm{peak}}=989$ TFLOPS，HBM3 峰值带宽 $B_w=3350$ GB/s。则 $I^*=$ \blank{2.5cm} FLOP/Byte，保留整数。若某 Decode 阶段算子的实测算术强度 $I=2$ FLOP/Byte，则该算子属于 \blank{3cm}，填写 Compute-bound 或 Memory-bound。}{}

\qitem{14}{对 GEMM $C_{M\times N}=A_{M\times K}B_{K\times N}$，若采用朴素实现，即计算每个 $C_{ij}$ 时从 HBM 重新读取 $A$ 的第 $i$ 行和 $B$ 的第 $j$ 列，$A,B$ 元素无任何复用，矩阵以 fp32 存储。则从 HBM 读取 $A,B$ 并写回 $C$ 的总数据量为 \blank{5cm} 字节。高效实现如 CUTLASS 通过 Tiling 使每个元素被复用多次，可将 HBM 访问量降低约 \blank{2.5cm} 倍，填数量级整数，如 10 或 100。}{}

\qitem{15}{CUDA Occupancy 指 SM 上活跃 Warp 数占最大支持 Warp 数的比例，受寄存器、Shared Memory 和 Block 大小限制。某 Kernel 每个 Thread Block 配置 256 个线程，即 8 个 Warp，每线程使用 64 个 32-bit 寄存器，不使用 Shared Memory。H100 每个 SM 的寄存器文件总量为 65536 个 32-bit 寄存器。则每个 SM 最多可同时驻留 \blank{2cm} 个该 Kernel 的 Thread Block。若 H100 每个 SM 最多支持 64 个 Warp，则该 Kernel 的理论最大 Occupancy 为 \blank{2.5cm}，填百分比如 $50\%$。}{}

\qitem{16}{SmoothQuant 对线性层 $\mathbf{Y}=\mathbf{X}\mathbf{W}$ 引入对角缩放矩阵 $S=\mathrm{diag}(s_1,\ldots,s_{d_{\mathrm{in}}})$，将激活离群点迁移到权重侧。数学等价的变换为 $\mathbf{Y}=$ \blank{5cm}。缩放因子公式为 $s_j=\dfrac{c_j^{\alpha}}{w_j^{1-\alpha}}$，其中 $c_j$ 为激活通道 $j$ 的最大绝对值，$w_j$ 为权重列 $j$ 的最大绝对值，$\alpha\in(0,1]$ 控制迁移强度。当 $\alpha=1$ 时，离群点完全由 \blank{2.5cm} 侧承担缩放，填写"激活"或"权重"。}{}

\qitem{17}{算子融合可减少中间张量的 HBM 往返。未融合时，RMSNorm 将输出 $\mathbf{h}\in\mathbb{R}^{d}$ 以 fp16 写回 HBM，$d$ 为隐藏维度，随后 Linear 层再从 HBM 读取 $\mathbf{h}$ 做 GEMM $\mathbf{Y}=\mathbf{h}W$。将 RMSNorm 与 Linear 融合为单个 Kernel 后，$\mathbf{h}$ 可保留在寄存器或 Shared Memory 中，节省 \blank{2cm} 次对 $\mathbf{h}$ 的完整 HBM 读写，其中 1 次写加 1 次读合为 2 次往返。若 Batch Size 为 $b$，则每前向步节省的 HBM 字节数为 \blank{3.5cm}。}{}

\qitem{18}{模型 FLOPs 利用率 MFU 是衡量大模型训练或推理硬件效率的核心指标，定义为实际达到的 FLOPs/s 与硬件峰值 FLOPs/s 之比：$\mathrm{MFU}=$ \blank{5cm}。PaLM 论文报告 540B 模型在 TPU v4 上 MFU 约 46\%。若 H100 峰值为 989 TFLOPS，实测某 70B 模型推理吞吐等价于 120 TFLOPS，则 MFU $=$ \blank{2.5cm}，保留整数百分比。}{}

\examsection{三、解答题}{本大题共 6 小题，共 60 分。解答应写出文字说明、证明过程或演算步骤。}

\qitem{19}{（本小题满分 10 分）}{}
\noindent 下图摘自 Williams et al., ``Roofline: An Insightful Visual Performance Model for Multicore Architectures'', CACM 2009, Figure 4，展示了算术强度与可达性能的关系。

\figplaceholder{请从 Roofline 论文 Figure 4 复制原图粘贴于此。该图展示了 Ridge Point 拐点及 Memory-bound 与 Compute-bound 两个性能区域。}{Williams et al., CACM 2009, Figure 4}

\noindent\textbf{已知：}算术强度 $I=\dfrac{\text{FLOPs}}{\text{HBM 访问字节数}}$；Ridge Point $I^*=\dfrac{P_{\mathrm{peak}}}{B_w}$；实际吞吐 $\approx\min\!\left(P_{\mathrm{peak}},\, I\cdot B_w\right)$。H100-80GB 参数：$P_{\mathrm{peak}}=989$ TFLOPS，$B_w=3350$ GB/s。

\begin{enumerate}[label=(\arabic*),leftmargin=2em,itemsep=0.2em,parsep=0pt,topsep=0pt]
  \item 结合上图，说明横轴、纵轴分别表示什么；标出 Ridge Point 的位置，并解释 $I<I^*$ 与 $I>I^*$ 时算子分别受何种资源限制；（4 分）
  \item 利用已知公式计算 H100 的 Ridge Point $I^*$，保留整数；（3 分）
  \item 某 GEMM 实测算术强度 $I=10$ FLOP/Byte，请用 Roofline 公式估算其可达吞吐，单位 TFLOPS，保留整数，并判断该算子属于 Memory-bound 还是 Compute-bound。（3 分）
\end{enumerate}

\answerlines{8}

\needspace{6cm}
\qitem{20}{（本小题满分 10 分）}{}
\noindent 自回归 Decode 阶段的性能瓶颈往往来自 KV Cache 的 HBM 带宽，而非算力不足。

\noindent\textbf{已知：}单条请求 KV Cache 显存，fp16 存储，为
\[
M_{\mathrm{KV}} = 2 \times L \times n \times h_{\mathrm{kv}} \times d_h \times 2\;\text{字节},
\]
其中 $L$ 为层数，$n$ 为当前序列长度，$h_{\mathrm{kv}}$ 为 KV 头数，$d_h$ 为每头维度，因子 $2$ 表示 K 与 V 各一份。Decode 每步需从 HBM 读取上述全部 KV，不重新计算历史 $K,V$。

\noindent 模型参数：$L=80$，$h_{\mathrm{kv}}=8$，$d_h=128$，$n=4096$，fp16。

\begin{enumerate}[label=(\arabic*),leftmargin=2em,itemsep=0.2em,parsep=0pt,topsep=0pt]
  \item 代入公式计算该请求 KV Cache 总显存 $M_{\mathrm{KV}}$，用 GB 表示，保留两位小数；（4 分）
  \item 估算 Decode 每步需从 HBM 读取的 KV 数据量，单位 GB，保留两位小数；（3 分）
  \item 结合 H100 Ridge Point $I^*\approx 295$ FLOP/Byte，用一句话说明为何 Decode 通常是 Memory-bound。（3 分）
\end{enumerate}

\answerlines{8}

\needspace{6cm}
\qitem{21}{（本小题满分 10 分）}{}
\noindent FlashAttention 通过 Tiling 在 SRAM 中完成分块注意力计算，避免将完整注意力矩阵写回 HBM。

\noindent\textbf{已知：}标准注意力分三步：$S=QK^{\top}$，$P=\mathrm{softmax}(S)$，$O=PV$。FlashAttention 的 Online Softmax 在遇到第 $j$ 个 $K,V$ 块时，用以下公式增量更新，其中 $m$ 为行最大值，$\ell$ 为归一化因子，$O$ 为输出累加器：
\[
m'=\max(m,\; m^{(j)}),\quad
\ell' = e^{m-m'}\ell + e^{m^{(j)}-m'}\ell^{(j)},\quad
O' = \frac{\ell}{\ell'}O\cdot e^{m-m'} + \frac{e^{m^{(j)}-m'}}{\ell'}O^{(j)}.
\]
HBM IO 复杂度：标准 Attention 写出 $S,P$ 为 $O(N^2)$；FlashAttention 为 $O(N^2d^2/M_{\mathrm{sram}})$。

\begin{enumerate}[label=(\arabic*),leftmargin=2em,itemsep=0.2em,parsep=0pt,topsep=0pt]
  \item 说明标准 Attention 将 $S,P$ 写回 HBM 在长序列下会带来什么问题；FlashAttention 如何避免；（3 分）
  \item 解释 $m'$ 更新公式中，为何需要比较旧最大值 $m$ 与新块最大值 $m^{(j)}$；$e^{m-m'}$ 项起什么作用；（4 分）
  \item 设 $N=8192$，$d=128$，$M_{\mathrm{sram}}=192$ KB，说明 FlashAttention 的 IO 复杂度相比 $O(N^2)$ 在长序列场景下的优势，无需精确计算数值。（3 分）
\end{enumerate}

\answerlines{10}

\needspace{6cm}
\qitem{22}{（本小题满分 10 分）}{}
\noindent\textbf{已知：}SmoothQuant 缩放因子 $s_j=\dfrac{c_j^{\alpha}}{w_j^{1-\alpha}}$，等价变换 $\mathbf{Y}=(\mathbf{X}S)(S^{-1}\mathbf{W})$，其中 $S=\mathrm{diag}(s_1,\ldots,s_d)$。GPTQ 量化第 $j$ 列权重后，按 Hessian 逆矩阵传播误差：
\[
\delta w_j = \hat{w}_j - w_j,\quad
w_i \leftarrow w_i - \frac{H^{-1}_{ij}}{H^{-1}_{jj}}\,\delta w_j \quad (i>j).
\]

\noindent 对某线性层第 3 个输入通道：$c_3=32$，$w_3=0.5$，$\alpha=0.5$。

\begin{enumerate}[label=(\arabic*),leftmargin=2em,itemsep=0.2em,parsep=0pt,topsep=0pt]
  \item 计算 $s_3$ 的值，保留两位小数，并说明 $\alpha=0.5$ 时激活与权重各承担多少量化难度；（3 分）
  \item 说明 $(\mathbf{X}S)(S^{-1}\mathbf{W})$ 为何与 $\mathbf{XW}$ 数学等价；从数值分布角度解释为何 $\mathbf{X}S$ 比 $\mathbf{X}$ 更易量化；（3 分）
  \item 设 GPTQ 量化 $w_3$ 后 $\hat{w}_3=0.48$，即 $\delta w_3=-0.02$，$H^{-1}_{33}=1.0$，$H^{-1}_{34}=0.1$，$w_4=1.0$。按上述公式计算更新后的 $w_4$。（4 分）
\end{enumerate}

\answerlines{8}

\needspace{6cm}
\qitem{23}{（本小题满分 10 分）}{}
\noindent 下图摘自 Kwon et al., ``Efficient Memory Management for Large Language Model Serving with PagedAttention'', SOSP 2023, Figure 4，展示了 PagedAttention 的 KV Cache 分块与映射机制。

\figplaceholder{请从 vLLM 论文 Figure 4 复制原图粘贴于此。该图展示了 KV Cache 划分为固定大小 Block，并通过 Block Table 映射到非连续物理页。}{Kwon et al., SOSP 2023, Figure 4}

\noindent 某团队计划在单卡 H100-80GB 上部署 70B 模型，fp16 权重约 140GB，需量化至 W8A8 约 70GB 方可放下，同时服务 100 个并发对话请求，平均上下文 4k tokens，峰值 32k tokens。模型：$L=80$，$h_{\mathrm{kv}}=8$，$d_h=128$，Block 大小 $B=16$ tokens。

\begin{enumerate}[label=(\arabic*),leftmargin=2em,itemsep=0.2em,parsep=0pt,topsep=0pt]
  \item 结合上图，解释 PagedAttention 如何解决 KV Cache 的内部碎片和外部碎片问题；（3 分）
  \item 计算单个请求在 $n=4096$ 时所需 KV Block 数量；100 个并发、平均 4k 时，KV Cache 总显存约多少 GB；（4 分）
  \item 当峰值 32k 请求与大量 4k 请求混合时，如何配置 Block 池大小与抢占策略以避免 OOM；（3 分）
\end{enumerate}

\answerlines{10}

\needspace{6cm}
\qitem{24}{（本小题满分 10 分，压轴）}{}
\noindent 某 LLM 推理服务部署在单卡 H100-80GB 上，70B 模型经 W8A8 量化后权重约 70GB，剩余约 10GB 可用于 KV Cache 与运行时缓冲。负载以 Decode 为主，平均上下文 2k tokens。系统观测到：SM 占用率 40\%；\texttt{dram\_throughput} 接近 90\% 峰值；\texttt{sm\_\_throughput} 仅 25\%。

\begin{enumerate}[label=(\arabic*),leftmargin=2em,itemsep=0.2em,parsep=0pt,topsep=0pt]
  \item 结合 Roofline 模型，$I^*\approx 295$ FLOP/Byte，分析上述 profiler 指标，说明瓶颈类型，并估算当前 Decode 算子的算术强度数量级；（3 分）
  \item 设计一套分层优化栈，至少从算子融合、W8A8 或 FP8 KV、FlashAttention-2、Continuous Batching 中选取 3 项，画出模块层次图或数据流，并说明每项在 Roofline 图上如何移动工作点；（4 分）
  \item 给出验收指标：优化后应观测到哪 3 个 Nsight 指标的变化，各朝什么方向变化即可认为优化生效。（3 分）
\end{enumerate}

\answerlines{12}

\newpage
\begin{center}
{\zihao{2}\heiti 参考答案与解析}\\[0.3cm]
{\small 命题：大模型高性能计算课程组 \quad 校对：教研组}
\end{center}
\vspace{0.3cm}
\noindent\rule{\linewidth}{0.6pt}

\subsection*{\heiti 一、选择题答案}

\begin{multicols}{4}
\begin{enumerate}[label=\textbf{\arabic*.},leftmargin=*,itemsep=0.1em,parsep=0pt,topsep=0pt]
  \item B
  \item A
  \item B
  \item B
  \item B
  \item C
  \item B
  \item B
  \item A
  \item B
  \item D
  \item B
\end{enumerate}
\end{multicols}

\vspace{0.3cm}
\noindent\textbf{1. B}\quad SM 内含 CUDA Core、Tensor Core、Warp 调度器、Shared Memory 和 Register File。Warp 32 线程是 SIMT 调度基本单位；分支发散会导致 Warp 内线程部分空闲。故选 B。

\noindent\textbf{2. A}\quad 越靠近计算单元，带宽越高、延迟越低：Register $>$ Shared Memory $>$ L2 $>$ HBM。故选 A。

\noindent\textbf{3. B}\quad $I<I^*$ 时性能受 HBM 带宽限制，为 Memory-bound。故选 B。

\noindent\textbf{4. B}\quad $MNK$ 次 MAC，每次 2 FLOPs，合计 $2MNK$。故选 B。

\noindent\textbf{5. B}\quad TPU 脉动阵列数据流规整；GPU 为 SIMT 通用并行。故选 B。

\noindent\textbf{6. C}\quad 融合受寄存器压力、算子形状和数据依赖约束，不能保证任意两算子时间减半。故选 C。

\noindent\textbf{7. B}\quad W8A8 标准流程：INT8 GEMM、INT32 累加、再乘 scale 反量化。故选 B。

\noindent\textbf{8. B}\quad FP16 MMA 为 $m16n8k16$，累加器 FP32；仍需 Tiling。故选 B。

\noindent\textbf{9. A}\quad GPTQ 用 Hessian 逐列、AWQ 保护显著通道、SmoothQuant 平滑激活离群点，均为 PTQ。故选 A。

\noindent\textbf{10. B}\quad Triton 以 Block 为抽象，编译器处理 Thread/Warp 与合并访存。故选 B。

\noindent\textbf{11. D}\quad FlashAttention-2 不会将 IO 升到 $O(N^3)$。故选 D。

\noindent\textbf{12. B}\quad TVM/Ansor 自动搜索超参，支持 GPU 等多后端。故选 B。

\subsection*{\heiti 二、填空题答案}

\noindent\textbf{13.}\quad $295$；Memory-bound\quad $I^*=\dfrac{989\times10^{12}}{3350\times10^9}\approx295$ FLOP/Byte。$I=2\ll 295$，带宽受限。

\noindent\textbf{14.}\quad $4(MK+KN+MN)$；$100$ 或 $128$\quad 朴素实现无复用；Tiling 典型降低约 100 倍 HBM 访问。

\noindent\textbf{15.}\quad $4$；$50\%$\quad 每 Block 耗 $256\times64=16384$ 寄存器，$65536/16384=4$ Block；$4\times8=32$ Warp，占最大 64 Warp 的 50\%。

\noindent\textbf{16.}\quad $(\mathbf{X}S)(S^{-1}\mathbf{W})$；权重\quad $\alpha=1$ 时 $s_j=c_j$，缩放完全施加在权重侧。

\noindent\textbf{17.}\quad $2$；$4bd$ 字节\quad 节省 1 写 + 1 读；每元素 fp16 占 2 字节，共 $b\times d$ 元素，$2\times b\times d\times 2=4bd$。

\noindent\textbf{18.}\quad $\dfrac{\text{实际 achieved FLOPs/s}}{\text{硬件峰值 FLOPs/s}}$；$12\%$\quad $120/989\approx12.1\%$。

\subsection*{\heiti 三、解答题答案}

\subsubsection*{19. Roofline 模型图示与计算（10 分）}

\noindent\textbf{(1) 图示理解（4 分）}

横轴为算术强度 $I$，纵轴为可达性能。图中有斜线 $I\cdot B_w$ 与水平线 $P_{\mathrm{peak}}$，交点为 Ridge Point $I^*$。$I<I^*$ 时位于斜线区域，Memory-bound；$I>I^*$ 时位于平顶区域，Compute-bound。

\noindent\textbf{(2) Ridge Point（3 分）}

$I^*=\dfrac{989\times10^{12}}{3350\times10^9}\approx295$ FLOP/Byte。

\noindent\textbf{(3) 吞吐估算（3 分）}

$I=10\ll 295$，Memory-bound。可达吞吐 $\approx I\cdot B_w=10\times3350$ GFLOPS $=33.5$ TFLOPS，保留整数约 \textbf{34 TFLOPS}。

\subsubsection*{20. KV Cache 与 Decode 带宽（10 分）}

\noindent\textbf{(1) $M_{\mathrm{KV}}$（4 分）}

$M_{\mathrm{KV}}=2\times80\times4096\times8\times128\times2=2^{31}$ 字节 $=2.00$ GB。

\noindent\textbf{(2) 每步读取（3 分）}

每步读取全部 KV，约 \textbf{2.00 GB}。

\noindent\textbf{(3) Memory-bound 原因（3 分）}

Decode 每步 FLOPs 仅 $O(d)$，但需读取数 GB 级 KV，算术强度约 0.01--1 FLOP/Byte，远低于 $I^*=295$，故为 Memory-bound。

\subsubsection*{21. FlashAttention 公式理解（10 分）}

\noindent\textbf{(1) 标准 Attention 的问题（3 分）}

$N\times N$ 的 $S,P$ 写回 HBM 导致 IO 量 $O(N^2)$，长序列时带宽成为瓶颈。FlashAttention 在 SRAM 分块计算，$S,P$ 不写出。

\noindent\textbf{(2) $m'$ 与 rescale（4 分）}

需比较新旧最大值以保证全局 softmax 正确；当 $m$ 更新为更大的 $m'$ 时，旧块贡献须按 $e^{m-m'}$ 重新缩放。

\noindent\textbf{(3) 复杂度优势（3 分）}

标准 Attention IO 至少 $O(N^2)$；FlashAttention 为 $O(N^2d^2/M_{\mathrm{sram}})$，$M_{\mathrm{sram}}$ 为常数时随 $N$ 增长更慢，长序列优势显著。

\subsubsection*{22. 量化公式应用（10 分）}

\noindent\textbf{(1) $s_3$（3 分）}

$s_3=\dfrac{\sqrt{32}}{\sqrt{0.5}}=\sqrt{64}=8.00$。$\alpha=0.5$ 时激活与权重各承担一半量化难度。

\noindent\textbf{(2) 等价性与易量化（3 分）}

$(\mathbf{X}S)(S^{-1}\mathbf{W})=\mathbf{XW}$。$\mathbf{X}S$ 压缩离群点动态范围，INT8 量化误差降低。

\noindent\textbf{(3) GPTQ 更新（4 分）}

$w_4\leftarrow 1.0-\dfrac{0.1}{1.0}\times(-0.02)=1.002$。

\subsubsection*{23. PagedAttention 架构设计（10 分）}

\noindent\textbf{(1) 碎片问题（3 分）}

固定 Block 分配消除内部碎片；Block Table 映射逻辑连续 KV 到非连续物理 Block，释放的 Block 可立即复用，消除外部碎片。

\noindent\textbf{(2) Block 数量与总显存（4 分）}

单请求 Block 数 $=\lceil 4096/16\rceil=256$。单请求 KV 约 2 GB，100 并发约 200 GB，需多卡或 KV 量化/卸载。

\noindent\textbf{(3) 调度策略（3 分）}

预分配 Block 池；长请求限制并发或 KV Offload；短请求优先；OOM 时按 LRU 抢占低优先级请求。

\subsubsection*{24. H100 综合架构设计（10 分）}

\noindent\textbf{(1) Roofline 分析（3 分）}

\texttt{dram\_throughput} 高、\texttt{sm\_\_throughput} 低，Memory-bound。Decode 算术强度约 0.1--2 FLOP/Byte，远低于 $I^*=295$。

\noindent\textbf{(2) 优化栈（4 分）}

W8A8/FP8 KV 减半 HBM 读取；Kernel Fusion 提高有效 $I$；Continuous Batching 增大有效 Batch；FlashAttention-2 优化 attention 带宽利用。

\noindent\textbf{(3) 验收指标（3 分）}

\texttt{sm\_\_throughput} 应上升；\texttt{dram\_bytes\_read}/token 应下降；\texttt{achieved\_occupancy} 应提升。

\vspace{0.5cm}
\noindent{\heiti 配图索引：}
\begin{itemize}[leftmargin=2em,itemsep=0.1em,parsep=0pt,topsep=0pt]
  \item \textbf{第 19 题}：Williams et al., ``Roofline Model'', CACM 2009, Figure 4
  \item \textbf{第 23 题}：Kwon et al., ``PagedAttention / vLLM'', SOSP 2023, Figure 4
\end{itemize}

\end{document}
